\subsection{Finiteness of extension and torsion modules}

Let M be a left A-module, I a filtered preordered set, \((\mathbf{N}_{i}, u_{ji})\) an inductive system of left A-modules relative to I, \(N = \varinjlim N_{i}\) its inductive limit, \(u_{i}: N_{i} \to N, i \in I\) , the canonical map. Then \((\operatorname{Ext}_{\mathbf{A}}(\mathbf{M}, \mathbf{N}_{i}), \operatorname{Ext}_{\mathbf{A}}(1_{\mathbf{M}}, u_{ji}))\) is an inductive system of k-modules and \((\operatorname{Ext}_{\mathbf{A}}(1_{\mathbf{M}}, u_{i}))\) an inductive system of maps, whose inductive limit is a homomorphism of graded k-modules, called canonical

\[
\varinjlim_ {i \in I} \operatorname{Ext} _ {\mathbf {A}} (\mathbf {M}, \mathbf {N} _ {i}) \rightarrow \operatorname{Ext} _ {\mathbf {A}} \left(\mathbf {M}, \varinjlim_ {i \in I} \mathbf {N} _ {i}\right).\tag{8}
\]

PROPOSITION 5. --- If A is left Noetherian and M is a finitely generated A-module, the canonical homomorphism (8) is bijective.

Indeed, let (X, p.~53, prop. 6) \(p: L \to M\) be a free resolution of M such that \(L_{n}\) be finitely generated for each \(n\) . The canonical morphism from \(k\) -complexes

\[
u: \varinjlim \operatorname{Homgr} _ {\mathbf {A}} (\mathrm{L}, \mathrm{N} _ {i}) \rightarrow \operatorname{Homgr} _ {\mathbf {A}} (\mathrm{L}, \varinjlim \mathrm{N} _ {i})
\]

is bijective, hence so is the homomorphism

\[
\varinjlim \mathrm{H} (\operatorname{Homgr} _ {\mathrm{A}} (\mathrm{L}, \mathrm{N} _ {i})) \rightarrow \mathrm{H} (\operatorname{Homgr} _ {\mathrm{A}} (\mathrm{L}, \varinjlim \mathrm{N} _ {i}))
\]

deduces homomorphisms \(\mathrm{H}(\mathrm{Homgr}(1,u_{i}))\) (X, p.~28, prop. 1). We then conclude by prop. 2 (X, p.~103) and th. 1 (X, p.~100).

PROPOSITION 6. --- Let B be a ring and N an (A, B)-bimodule, which is a Noetherian B-module (resp. of finite length).

\begin{enumerate}
\def\labelenumi{\alph{enumi})}
\item
  Suppose A is right Noetherian and let M be a finitely generated right A-module. Then the B-modules (X, p.~81) Tor \(_{n}^{A}\) (M, N) are Noetherian (resp. of finite length).
\item
  Suppose A is left Noetherian and let M be a finitely generated left A-module. Then the B-modules (X, p.~98) Ext \(_{A}^{n}\) (M, N) are Noetherian (resp. of finite length).
\end{enumerate}

Choose a free resolution \(p: \mathbf{L} \to \mathbf{M}\) such that each of the A-modules \(\mathbf{L}_n\) be finitely generated (X, p.~53, prop. 6), and let C be the complex of B-modules \(\mathbf{L} \otimes_{\mathbf{A}} \mathbf{N}\) in the case \(a\) ), \(\operatorname{Homgr}_{\mathbf{A}}(\mathbf{L}, \mathbf{N})\) in the case \(b\) ). Each of the B-modules \(\mathbf{C}_n\) is isomorphic to a product of a finite number of copies of N, hence is Noetherian (resp. of finite length); the same is therefore true of the modules \(\mathrm{H}_n(\mathbf{C})\) . Now, according to X, p.~100, th. 1, these are isomorphic to the \(\operatorname{Tor}_n^{\mathbf{A}}(\mathbf{M}, \mathbf{N})\) in the case \(a\) ), to the \(\operatorname{Ext}_{\mathbf{A}}^{-n}(\mathbf{M}, \mathbf{N})\) in the case \(b\) ).

COROLLARY. --- Let \(\rho: A \to B\) be a homomorphism of Noetherian commutative rings, M a finitely generated A-module, N a B-module. If N is a finitely generated (resp. finite length) B-module, then so are the B-modules \(\operatorname{Tor}_{n}^{A}(M, N)\) and \(\operatorname{Ext}_{A}^{n}(M, N)\) .

